\subsection{标量环的限制与扩张之间的关系}

设 \(\rho: \mathbf{A} \to \mathbf{B}\) 为环同态。对每个左A-模E，典范A-线性映射

\[
\phi_ {\mathbf {E}}: \mathbf {E} \rightarrow \rho_ {*} (\rho^ {*} (\mathbf {E}))\tag{10}
\]

已在第1号中定义，使得 \(\phi_{\mathbb{E}}(x) = 1 \otimes x\)。现在我们考虑左B-模F，并将命题1（第1号）应用于A-同态 \(\mathbf{l}_{\rho_{*}(\mathbf{F})}: \rho_{*}(\mathbf{F}) \to \rho_{*}(\mathbf{F})\)：我们得到一个B-线性映射

\[
\psi_ {\mathbf {F}} \colon \rho^ {*} (\rho_ {*} (\mathbf {F})) \rightarrow \mathbf {F}\tag{11}
\]

等于 \(\delta(1_{\rho_*(\mathbb{F})})\) 并且因此使得，对所有 \(y \in \mathbb{F}\) 以及所有 \(\beta \in \mathbb{B}\) , \(\psi_{\mathbb{F}}(\beta \otimes y) = \beta y\) .

命题5. 设E为左A-模，F为左B-模；复合映射

\[
\rho^ {*} (\mathrm{E}) \xrightarrow [ \rho^ {*} (\phi_ {\mathrm{E}}) ]{} \rho^ {*} (\rho_ {*} (\rho^ {*} (\mathrm{E}))) \xrightarrow [ \psi_ {\rho^ {*} (\mathrm{E})} ]{} \rho^ {*} (\mathrm{E})\tag{12}
\]

\[
\rho_ {*} (\mathrm{F}) \xrightarrow [ \phi_ {\rho_ {*} (\mathrm{F})} ]{} \rho_ {*} (\rho^ {*} (\rho_ {*} (\mathrm{F}))) \xrightarrow [ \rho_ {*} (\psi_ {\mathrm{F}}) ]{} \rho_ {*} (\mathrm{F})\tag{13}
\]

分别等于 \(\rho^{*}(\mathbf{E})\) 与 \(\rho_{*}(\mathbf{F})\) 的恒等映射。

我们给出(12)的证明作为例子；对所有 \(x \in \mathbf{E}\)，映射 \(\rho^{*}(\phi_{\mathbf{E}})\) 把 \(1 \otimes x\) 映为元素 \(1 \otimes (1 \otimes x)\)，且映射 \(\psi_{\rho^{*}(\mathbf{E})}\) 把 \(1 \otimes (1 \otimes x)\) 映为元素 \(1 \otimes x\)；结论由形如 \(1 \otimes x\) 的元素生成B-模 \(\rho^{*}(\mathbf{E})\) 这一事实得出；(13)的证明甚至更简单。

推论. 映射 \(\rho^{*}(\phi_{\mathbf{E}})\) 和 \(\phi_{\rho^{*}(\mathbf{F})}\) 是单射，并且分别把 \(\rho^{*}(\mathbf{E})\) 等同于 \(\rho^{*}(\rho_{*}(\rho^{*}(\mathbf{E})))\) 的一个直因子，把 \(\rho_{*}(\mathbf{F})\) 等同于 \(\rho_{*}(\rho^{*}(\rho_{*}(\mathbf{F})))\) 的一个直因子。

这是命题5及§1第9号中命题15的推论2的直接结果。

命题6. 设E为左A-模，F为右B-模。则存在唯一的Z-同态

\[
\rho_ {*} (\mathrm{F}) \otimes_ {\mathrm{A}} \mathrm{E} \rightarrow \mathrm{F} \otimes_ {\mathrm{B}} \rho^ {*} (\mathrm{E})\tag{14}
\]

它把 \(y \otimes x\) 映为 \(y \otimes (1 \otimes x)\)（对所有 \(x \in E\) 及所有 \(y \in F\)），且该同态是双射。

根据定义，（14）的右边是 \(F \otimes_{B} (B \otimes_{A} E)\)，其中B被视为(B, A)-双模，并且存在一个§3第8号命题8中定义的典范Z-同构 \((\mathbf{F} \otimes_{\mathbf{B}} \mathbf{B}) \otimes_{\mathbf{A}} \mathbf{E} \to \mathbf{F} \otimes_{\mathbf{B}} (\mathbf{B} \otimes_{\mathbf{A}} \mathbf{E})\)；另一方面，§3第4号命题4中的典范同构 \(F \to F \otimes_{B} B\) 对由 \(\rho\) 定义的两侧右A-模结构是同构。由此得到所期望的同构。

当A与B交换时，同构(14)是A-模同构

\[
\rho_ {*} (\mathrm{F}) \otimes_ {\mathrm{A}} \mathrm{E} \rightarrow \rho_ {*} (\mathrm{F} \otimes_ {\mathrm{B}} \rho^ {*} (\mathrm{E})).
\]

